% Claim proof file for row claim_3_17 -- "MarginalizationsCommute".
% Auto-generated stub by scaffold/solve_chapter.py at row start. A manager
% agent dedicated to writing the TeX proof fills in the proof body below.
%
% This file is a sibling of `main.tex` inside `<subsection>/tex/`. The
% `subfiles` package lets it render both standalone and as part of
% `main.tex` via `\subfile{...}`.
%
% The statement is restated above the proof so the file is self-contained
% when rendered on its own. The proof must:
%   - Stay close to the lecture notes' proof if one exists (always search
%     the LN first for a `\begin{proof}` block immediately following the
%     claim; copy / lightly edit when present).
%   - Otherwise use the LN paradigm and cite prior claims by their `ref`.
%   - Be self-contained enough that a mathematician can follow it without
%     re-reading the LN.
%   - Have no `sorry`, `omitted`, or "obvious" shortcuts.

\documentclass[main]{subfiles}

\begin{document}
% ref: claim_3_17 (proof, with statement restated for readability)
% title: MarginalizationsCommute
\def\rowref{claim\_3\_17}
\phantomsection\label{claim_3_17}

% --- statement (restated) ---------------------------------------------------
\begin{Lem}[Marginalizations commute]\label{marginalizations-commute}
      Let $G=(J,V,E,L)$ be a CDMG and $W_1, W_2 \ins V$ two disjoint subsets of output nodes.
      Then we have:
      \[ \lp G^{\sm W_1} \rp^{\sm W_2} = \lp G^{\sm W_2} \rp^{\sm W_1} =  G^{\sm (W_1 \cup W_2)}. \]
     % This shows that the notation:
     % \[ G(V \sm (W_1 \cup W_2)|\doit(J))  \]
     % is unambiguous.
\end{Lem}

% --- proof ------------------------------------------------------------------
\begin{proof}
% TODO: write the proof body.
\end{proof}

\end{document}
